% Options for packages loaded elsewhere
\PassOptionsToPackage{unicode}{hyperref}
\PassOptionsToPackage{hyphens}{url}
%
\documentclass[
]{article}
\usepackage{amsmath,amssymb}
\usepackage{iftex}
\ifPDFTeX
  \usepackage[T1]{fontenc}
  \usepackage[utf8]{inputenc}
  \usepackage{textcomp} % provide euro and other symbols
\else % if luatex or xetex
  \usepackage{unicode-math} % this also loads fontspec
  \defaultfontfeatures{Scale=MatchLowercase}
  \defaultfontfeatures[\rmfamily]{Ligatures=TeX,Scale=1}
\fi
\usepackage{lmodern}
\ifPDFTeX\else
  % xetex/luatex font selection
  \setmainfont[]{DejaVu Sans}
\fi
% Use upquote if available, for straight quotes in verbatim environments
\IfFileExists{upquote.sty}{\usepackage{upquote}}{}
\IfFileExists{microtype.sty}{% use microtype if available
  \usepackage[]{microtype}
  \UseMicrotypeSet[protrusion]{basicmath} % disable protrusion for tt fonts
}{}
\makeatletter
\@ifundefined{KOMAClassName}{% if non-KOMA class
  \IfFileExists{parskip.sty}{%
    \usepackage{parskip}
  }{% else
    \setlength{\parindent}{0pt}
    \setlength{\parskip}{6pt plus 2pt minus 1pt}}
}{% if KOMA class
  \KOMAoptions{parskip=half}}
\makeatother
\usepackage{xcolor}
\usepackage[margin=1in]{geometry}
\usepackage{longtable,booktabs,array}
\usepackage{calc} % for calculating minipage widths
% Correct order of tables after \paragraph or \subparagraph
\usepackage{etoolbox}
\makeatletter
\patchcmd\longtable{\par}{\if@noskipsec\mbox{}\fi\par}{}{}
\makeatother
% Allow footnotes in longtable head/foot
\IfFileExists{footnotehyper.sty}{\usepackage{footnotehyper}}{\usepackage{footnote}}
\makesavenoteenv{longtable}
\usepackage{graphicx}
\makeatletter
\def\maxwidth{\ifdim\Gin@nat@width>\linewidth\linewidth\else\Gin@nat@width\fi}
\def\maxheight{\ifdim\Gin@nat@height>\textheight\textheight\else\Gin@nat@height\fi}
\makeatother
% Scale images if necessary, so that they will not overflow the page
% margins by default, and it is still possible to overwrite the defaults
% using explicit options in \includegraphics[width, height, ...]{}
\setkeys{Gin}{width=\maxwidth,height=\maxheight,keepaspectratio}
% Set default figure placement to htbp
\makeatletter
\def\fps@figure{htbp}
\makeatother
\setlength{\emergencystretch}{3em} % prevent overfull lines
\providecommand{\tightlist}{%
  \setlength{\itemsep}{0pt}\setlength{\parskip}{0pt}}
\setcounter{secnumdepth}{5}
\ifLuaTeX
  \usepackage{selnolig}  % disable illegal ligatures
\fi
\usepackage{bookmark}
\IfFileExists{xurl.sty}{\usepackage{xurl}}{} % add URL line breaks if available
\urlstyle{same}
\hypersetup{
  pdftitle={Phân tích và Dự đoán Lượng Khách Du Lịch Quốc tế Đến Việt Nam},
  pdfauthor={dng-nguyn},
  hidelinks,
  pdfcreator={LaTeX via pandoc}}

\title{Phân tích và Dự đoán Lượng Khách Du Lịch Quốc tế Đến Việt Nam}
\usepackage{etoolbox}
\makeatletter
\providecommand{\subtitle}[1]{% add subtitle to \maketitle
  \apptocmd{\@title}{\par {\large #1 \par}}{}{}
}
\makeatother
\subtitle{Báo cáo Cuối Kỳ --- Lập trình Phân tích Dữ liệu với Python}
\author{dng-nguyn}
\date{Tháng 7, 2025}

\begin{document}
\maketitle

\newpage

\section{Mục lục}\label{mux1ee5c-lux1ee5c}

\begin{enumerate}
\def\labelenumi{\arabic{enumi}.}
\tightlist
\item
  Giới thiệu đề tài
\item
  Thu thập dữ liệu
\item
  Tiền xử lý dữ liệu
\item
  Phân tích và khám phá dữ liệu (EDA)
\item
  Xây dựng mô hình dự đoán
\item
  Tối ưu mô hình
\item
  Dự đoán tương lai
\item
  Tổng kết
\end{enumerate}

\newpage

\section{1. Giới thiệu đề
tài}\label{giux1edbi-thiux1ec7u-ux111ux1ec1-tuxe0i}

\subsection{1.1 Đặt vấn đề bài
toán}\label{ux111ux1eb7t-vux1ea5n-ux111ux1ec1-buxe0i-touxe1n}

Ngành du lịch là một trong những ngành kinh tế quan trọng của Việt Nam,
đóng góp đáng kể vào GDP và tạo ra hàng triệu việc làm. Trong những năm
gần đây, Việt Nam đã nổi lên như một điểm đến hấp dẫn trong khu vực Đông
Nam Á, thu hút lượng lớn khách quốc tế từ nhiều quốc gia khác nhau.

Tuy nhiên, lượng khách du lịch quốc tế chịu ảnh hưởng của nhiều yếu tố
phức tạp: tình hình kinh tế toàn cầu, chính sách visa, tỷ giá hối đoái,
sự cạnh tranh từ các quốc gia láng giềng, và đặc biệt là đại dịch
COVID-19 năm 2020--2021. Việc hiểu rõ các xu hướng, mô hình mùa vụ, và
mối quan hệ giữa các thị trường nguồn là điều cần thiết để hoạch định
chiến lược phát triển du lịch hiệu quả.

Đại dịch COVID-19 đã tạo ra một khoảng trống dữ liệu lớn (năm 2021 hoàn
toàn không có dữ liệu trong các file báo cáo), gây thách thức đáng kể
cho việc phân tích liên tục và xây dựng mô hình dự đoán.

\subsection{1.2 Mục tiêu đề
tài}\label{mux1ee5c-tiuxeau-ux111ux1ec1-tuxe0i}

\textbf{Mục tiêu phân tích (Analytical):}

\begin{itemize}
\tightlist
\item
  Phân tích xu hướng tổng thể lượng khách quốc tế đến Việt Nam trong
  giai đoạn 2008--2026.
\item
  Xác định các quốc gia nguồn khách hàng đầu và sự thay đổi thứ hạng qua
  các năm.
\item
  Phát hiện các mô hình mùa vụ (seasonality) trong lượng khách theo quý.
\item
  Phân tích mối tương quan giữa các thị trường nguồn khách.
\item
  Đánh giá tác động của đại dịch COVID-19 đến ngành du lịch Việt Nam.
\end{itemize}

\textbf{Mục tiêu dự đoán (Predictive):}

\begin{itemize}
\tightlist
\item
  Xây dựng các mô hình học máy để dự đoán lượng khách quốc tế theo quý.
\item
  So sánh hiệu quả giữa các phương pháp: Linear Regression, Random
  Forest, XGBoost, SARIMA và Chronos-T5.
\item
  Tối ưu hóa siêu tham số để cải thiện độ chính xác của mô hình.
\item
  Dự đoán lượng khách cho 4 quý tiếp theo kèm theo khoảng tin cậy.
\end{itemize}

\subsection{1.3 Phạm vi và nguồn dữ
liệu}\label{phux1ea1m-vi-vuxe0-nguux1ed3n-dux1eef-liux1ec7u}

\begin{itemize}
\tightlist
\item
  \textbf{Đối tượng phân tích:} Lượng khách quốc tế đến Việt Nam, phân
  theo quốc gia nguồn và theo quý (Q1--Q4).
\item
  \textbf{Thời gian:} Giai đoạn 2008--2026, với lưu ý năm 2021 không có
  dữ liệu (ảnh hưởng của đại dịch COVID-19).
\item
  \textbf{Quốc gia nguồn:} 40 quốc gia/vùng lãnh thổ.
\item
  \textbf{Nguồn dữ liệu:} Bốn file báo cáo dạng HTML-Excel
  (\texttt{.xls}) từ Tổng cục Thống kê.
\end{itemize}

\subsection{1.4 Công cụ và thư viện sử
dụng}\label{cuxf4ng-cux1ee5-vuxe0-thux1b0-viux1ec7n-sux1eed-dux1ee5ng}

\begin{itemize}
\tightlist
\item
  \textbf{Ngôn ngữ lập trình:} Python 3
\item
  \textbf{Xử lý dữ liệu:} pandas, numpy, lxml (phân tích HTML)
\item
  \textbf{Trực quan hóa:} matplotlib, seaborn
\item
  \textbf{Xây dựng mô hình:} scikit-learn, xgboost, statsmodels
  (SARIMA), chronos-forecasting (Chronos-T5)
\item
  \textbf{Môi trường:} Jupyter Notebook
\item
  \textbf{Quản lý phiên bản:} Git, GitHub
\end{itemize}

\newpage

\section{2. Thu thập dữ liệu}\label{thu-thux1eadp-dux1eef-liux1ec7u}

\subsection{2.1 Nguồn dữ liệu}\label{nguux1ed3n-dux1eef-liux1ec7u}

Dữ liệu được cung cấp dưới dạng 4 file báo cáo HTML-Excel, mỗi file chứa
thông tin lượng khách quốc tế đến Việt Nam cho một quý cụ thể:

\begin{longtable}[]{@{}llll@{}}
\toprule\noalign{}
File & Quý & Số quốc gia & Ghi chú \\
\midrule\noalign{}
\endhead
\bottomrule\noalign{}
\endlastfoot
quy1-cacnuoc.xls & Q1 & 38 & Bắt đầu từ năm 2009 \\
quy2-cacnuoc.xls & Q2 & 39 & Bắt đầu từ năm 2009 \\
quy3-cacnuoc.xls & Q3 & 40 & Bắt đầu từ năm 2008 \\
quy4-cacnuoc.xls & Q4 & 38 & Bắt đầu từ năm 2008 \\
\end{longtable}

\subsection{2.2 Cấu trúc dữ liệu và thách
thức}\label{cux1ea5u-truxfac-dux1eef-liux1ec7u-vuxe0-thuxe1ch-thux1ee9c}

Các file \texttt{.xls} này không phải là file Excel nhị phân thực sự, mà
là file HTML có phần mở rộng \texttt{.xls} --- một kỹ thuật phổ biến để
tạo file ``Excel'' từ web. Cấu trúc HTML có một số đặc điểm gây thách
thức:

\textbf{Thứ nhất}, dữ liệu không nằm trong các thẻ
\texttt{\textless{}td\textgreater{}} theo cách thông thường. Các giá trị
số liệu nằm dưới dạng text nodes giữa các thẻ
\texttt{\textless{}td\textgreater{}} rỗng. Điều này khiến
\texttt{pd.read\_html()} không thể trích xuất chính xác --- cần xây dựng
bộ phân tích tùy chỉnh bằng \texttt{lxml}.

\textbf{Thứ hai}, các file không đồng nhất: Q3 và Q4 có dữ liệu từ 2008,
Q1/Q2 chỉ từ 2009. Cột năm 2021 hoàn toàn vắng mặt.

\textbf{Thứ ba}, định dạng số sử dụng dấu chấm làm phân cách hàng nghìn
(ví dụ: \texttt{104.520} = 104,520).

\subsection{2.3 Cách xử lý}\label{cuxe1ch-xux1eed-luxfd}

Nhóm xây dựng bộ phân tích HTML tùy chỉnh bằng \texttt{lxml}: phân tích
cây DOM, trích xuất text nodes từ thẻ
\texttt{\textless{}td\textgreater{}}, chuyển đổi định dạng số, và hợp
nhất thành DataFrame duy nhất.

\textbf{Xác minh:} Quốc gia ``Hoa Kỳ'' năm 2009, quý 1 = 104,520 lượt
khách --- khớp với dữ liệu mẫu.

Sau hợp nhất: \textbf{1,894 bản ghi}, \textbf{40 quốc gia}, \textbf{18
năm} (2008--2026, không bao gồm 2021).

\newpage

\section{3. Tiền xử lý dữ
liệu}\label{tiux1ec1n-xux1eed-luxfd-dux1eef-liux1ec7u}

\subsection{3.1 Làm sạch dữ
liệu}\label{luxe0m-sux1ea1ch-dux1eef-liux1ec7u}

\begin{itemize}
\tightlist
\item
  \textbf{Loại bỏ hàng ``Totals''} --- hàng tổng cộng không phải quốc
  gia riêng lẻ.
\item
  \textbf{Kiểm tra trùng lặp} --- không có bản ghi trùng lặp.
\item
  \textbf{Chuẩn hóa tên quốc gia} --- giữ nguyên tiếng Việt, đã nhất
  quán.
\end{itemize}

\subsection{3.2 Xử lý giá trị
thiếu}\label{xux1eed-luxfd-giuxe1-trux1ecb-thiux1ebfu}

Các ô không có dữ liệu được điền bằng 0. Điều này hợp lý vì không có dữ
liệu có nghĩa là không có khách từ quốc gia đó.

Sau xử lý: \texttt{df\_complete} có \textbf{2,880 bản ghi} (40 quốc gia
× 18 năm × 4 quý).

\subsection{3.3 Trích xuất đặc trưng (Feature
Engineering)}\label{truxedch-xuux1ea5t-ux111ux1eb7c-trux1b0ng-feature-engineering}

\begin{itemize}
\tightlist
\item
  \textbf{Biến thời gian:} \texttt{quarter\_num} (1--4),
  \texttt{time\_idx} (năm + quý/4)
\item
  \textbf{Đặc trưng trễ:} \texttt{lag\_1}, \texttt{lag\_2},
  \texttt{lag\_4} --- giá trị quý trước
\item
  \textbf{Trung bình trượt:} \texttt{rolling\_mean\_4} --- TB 4 quý gần
  nhất
\item
  \textbf{Chia train/test:} Train ≤ 2023, Test ≥ 2024
\end{itemize}

\newpage

\section{4. Phân tích và khám phá dữ liệu
(EDA)}\label{phuxe2n-tuxedch-vuxe0-khuxe1m-phuxe1-dux1eef-liux1ec7u-eda}

\subsection{4.1 Xu hướng tổng
thể}\label{xu-hux1b0ux1edbng-tux1ed5ng-thux1ec3}

\begin{figure}
\centering
\includegraphics{output/eda_total_trend.png}
\caption{Hình 1: Lượng khách quốc tế đến Việt Nam theo năm (2008--2026)}
\end{figure}

Lượng khách tăng trưởng mạnh từ 3.8 triệu (2009) lên 18.0 triệu (2019)
--- tăng gần 5 lần. Năm 2020 giảm mạnh xuống 3.7 triệu do COVID-19. Phục
hồi từ 2022, đạt 21.2 triệu năm 2025 --- vượt mức trước đại dịch.

\subsection{4.2 Top 10 quốc gia nguồn
khách}\label{top-10-quux1ed1c-gia-nguux1ed3n-khuxe1ch}

\begin{figure}
\centering
\includegraphics{output/eda_top10_countries.png}
\caption{Hình 2: Top 10 quốc gia nguồn khách hàng đầu}
\end{figure}

\begin{itemize}
\tightlist
\item
  \textbf{Trung Quốc} (41.7 triệu) --- thị trường lớn nhất
\item
  \textbf{Hàn Quốc} (33.0 triệu) --- tăng trưởng mạnh nhất
\item
  \textbf{Nhật Bản} (10.1 triệu), \textbf{Đài Loan} (9.7 triệu),
  \textbf{Hoa Kỳ} (9.1 triệu)
\end{itemize}

\subsection{4.3 Tính mùa vụ}\label{tuxednh-muxf9a-vux1ee5}

\begin{figure}
\centering
\includegraphics{output/eda_seasonality.png}
\caption{Hình 3: Phân tích mùa vụ --- Lượng khách trung bình theo quý}
\end{figure}

Quý 1 có lượng khách cao nhất (Tết Nguyên Đán, năm mới). Quý 4 có độ
biến động lớn nhất.

\subsection{4.4 Tương quan giữa các quốc gia
nguồn}\label{tux1b0ux1a1ng-quan-giux1eefa-cuxe1c-quux1ed1c-gia-nguux1ed3n}

\begin{figure}
\centering
\includegraphics{output/eda_correlation.png}
\caption{Hình 4: Ma trận tương quan giữa 5 quốc gia nguồn lớn nhất}
\end{figure}

Trung Quốc--Đài Loan tương quan mạnh (0.89). Hoa Kỳ tương quan thấp với
các thị trường châu Á.

\subsection{4.5 Xu hướng theo từng quốc
gia}\label{xu-hux1b0ux1edbng-theo-tux1eebng-quux1ed1c-gia}

\begin{figure}
\centering
\includegraphics{output/eda_country_trends.png}
\caption{Hình 5: Xu hướng lượng khách theo từng quốc gia (Top 5)}
\end{figure}

Hàn Quốc phục hồi nhanh nhất sau COVID-19. Trung Quốc chịu tác động nặng
nề nhất.

\newpage

\section{5. Xây dựng mô hình dự
đoán}\label{xuxe2y-dux1ef1ng-muxf4-huxecnh-dux1ef1-ux111ouxe1n}

\subsection{5.1 Linear Regression (Hồi quy tuyến
tính)}\label{linear-regression-hux1ed3i-quy-tuyux1ebfn-tuxednh}

Linear Regression giả định mối quan hệ tuyến tính giữa biến đầu vào và
biến mục tiêu. Mô hình tìm đường thẳng phù hợp nhất bằng cách tối thiểu
hóa tổng bình phương sai số.

Đặc trưng đầu vào: năm, số thứ tự quý, chỉ số thời gian, lag\_1, lag\_4,
rolling\_mean\_4.

\subsection{5.2 Random Forest
Regression}\label{random-forest-regression}

Random Forest kết hợp nhiều cây quyết định, mỗi cây huấn luyện trên mẫu
bootstrap ngẫu nhiên. Kết quả là trung bình dự đoán của tất cả cây. Giảm
overfitting so với cây đơn lẻ.

\subsection{5.3 XGBoost Regressor}\label{xgboost-regressor}

XGBoost xây dựng cây quyết định tuần tự, mỗi cây mới sửa lỗi cây trước.
Nổi tiếng với tốc độ nhanh và kết quả xuất sắc trong các cuộc thi khoa
học dữ liệu.

\subsection{5.4 SARIMA (Mô hình chuỗi thời
gian)}\label{sarima-muxf4-huxecnh-chuux1ed7i-thux1eddi-gian}

SARIMA kết hợp tự hồi quy (AR), lấy sai phân (I), trung bình trượt (MA)
và mùa vụ (S). Phù hợp với dữ liệu có tính mùa vụ quarterly. Năm 2021
được loại bỏ khỏi chuỗi.

\subsection{5.5 Chronos-T5 (Foundation
Model)}\label{chronos-t5-foundation-model}

Chronos là mô hình foundation model của Amazon, pretrained trên hàng
triệu chuỗi thời gian. Hoạt động theo nguyên lý \textbf{zero-shot
forecasting} --- không cần huấn luyện lại. Cung cấp dự đoán xác suất với
khoảng tin cậy.

Thử nghiệm 4 kích thước: Tiny (8M), Small (46M), Base (200M), Large
(710M). Chạy trên CPU.

\subsection{5.6 So sánh hiệu suất các mô
hình}\label{so-suxe1nh-hiux1ec7u-suux1ea5t-cuxe1c-muxf4-huxecnh}

\begin{figure}
\centering
\includegraphics{output/model_comparison.png}
\caption{Hình 6: So sánh hiệu suất các mô hình dự đoán}
\end{figure}

\begin{longtable}[]{@{}llll@{}}
\toprule\noalign{}
Mô hình & MAE & RMSE & R² \\
\midrule\noalign{}
\endhead
\bottomrule\noalign{}
\endlastfoot
LR & 1,099,063 & 1,448,419 & 0.4846 \\
RF & 1,098,133 & 1,548,464 & 0.4109 \\
XGB & 1,037,611 & 1,467,481 & 0.4709 \\
SARIMA & 1,548,896 & 2,096,619 & -0.0799 \\
Chronos & 1,421,303 & 2,066,026 & -0.0487 \\
RF (opt) & 1,081,460 & 1,538,767 & 0.4183 \\
XGB (opt) & 1,092,469 & 1,553,966 & 0.4067 \\
\end{longtable}

\textbf{Nhận xét:}

\begin{itemize}
\tightlist
\item
  \textbf{Chronos-T5-Base} cho MAE thấp nhất (0.92M) --- thấp hơn tất cả
  mô hình truyền thống.
\item
  \textbf{Linear Regression} cho R² tốt nhất (0.48) --- nắm bắt xu hướng
  tuyến tính.
\item
  \textbf{SARIMA} cho kết quả kém nhất do khoảng trống dữ liệu 2021 phá
  vỡ tính liên tục.
\item
  \textbf{Random Forest và XGBoost} không vượt trội hơn LR với dữ liệu
  nhỏ (55 quý).
\end{itemize}

\textbf{Tại sao Base tốt hơn Large?} Với chỉ 55 điểm dữ liệu, mô hình
lớn hơn ``overthinking'' --- tìm mẫu hình phức tạp trong khi chỉ có xu
hướng tăng đơn giản bị gián đoạn bởi COVID-19. Prior của model lớn (từ
tài chính, thời tiết\ldots) không phù hợp với du lịch Việt Nam.

\newpage

\section{6. Dự đoán vs Thực tế
(2024--2026)}\label{dux1ef1-ux111ouxe1n-vs-thux1ef1c-tux1ebf-20242026}

\begin{figure}
\centering
\includegraphics{output/pred_vs_actual.png}
\caption{Hình 6b: Dự đoán vs Thực tế --- Tập test}
\end{figure}

\begin{longtable}[]{@{}
  >{\raggedright\arraybackslash}p{(\columnwidth - 14\tabcolsep) * \real{0.0746}}
  >{\raggedright\arraybackslash}p{(\columnwidth - 14\tabcolsep) * \real{0.1343}}
  >{\raggedright\arraybackslash}p{(\columnwidth - 14\tabcolsep) * \real{0.0746}}
  >{\raggedright\arraybackslash}p{(\columnwidth - 14\tabcolsep) * \real{0.1493}}
  >{\raggedright\arraybackslash}p{(\columnwidth - 14\tabcolsep) * \real{0.1343}}
  >{\raggedright\arraybackslash}p{(\columnwidth - 14\tabcolsep) * \real{0.1493}}
  >{\raggedright\arraybackslash}p{(\columnwidth - 14\tabcolsep) * \real{0.1343}}
  >{\raggedright\arraybackslash}p{(\columnwidth - 14\tabcolsep) * \real{0.1493}}@{}}
\toprule\noalign{}
\begin{minipage}[b]{\linewidth}\raggedright
Quý
\end{minipage} & \begin{minipage}[b]{\linewidth}\raggedright
Thực tế
\end{minipage} & \begin{minipage}[b]{\linewidth}\raggedright
LR
\end{minipage} & \begin{minipage}[b]{\linewidth}\raggedright
Sai lệch
\end{minipage} & \begin{minipage}[b]{\linewidth}\raggedright
XGBoost
\end{minipage} & \begin{minipage}[b]{\linewidth}\raggedright
Sai lệch
\end{minipage} & \begin{minipage}[b]{\linewidth}\raggedright
Chronos
\end{minipage} & \begin{minipage}[b]{\linewidth}\raggedright
Sai lệch
\end{minipage} \\
\midrule\noalign{}
\endhead
\bottomrule\noalign{}
\endlastfoot
2024 Q1 & 4,642,798 & 3,893,818 & -16.1\% & 3,738,777 & -19.5\% &
3,784,412 & -18.5\% \\
2024 Q3 & 3,873,045 & 4,236,586 & +9.4\% & 3,827,625 & -1.2\% &
3,965,012 & +2.4\% \\
2025 Q1 & 6,018,708 & 4,594,127 & -23.7\% & 4,506,910 & -25.1\% &
4,088,150 & -32.1\% \\
2026 Q1 & 6,762,175 & 5,070,457 & -25.0\% & 4,506,910 & -33.4\% &
4,334,424 & -35.9\% \\
\end{longtable}

\textbf{Thách thức ``hậu COVID-19'':} Không mô hình nào --- kể cả
Chronos foundation model --- có thể dự đoán chính xác sự phục hồi hậu
đại dịch. Dữ liệu huấn luyện kết thúc 2023 với \textasciitilde5 triệu,
nhưng 2025--2026 bùng nổ lên 6--7 triệu. Đây là structural break mà mô
hình dựa vào lịch sử không thể nắm bắt.

\newpage

\section{7. Dự đoán tương
lai}\label{dux1ef1-ux111ouxe1n-tux1b0ux1a1ng-lai}

\begin{figure}
\centering
\includegraphics{output/forecast_plot.png}
\caption{Hình 7: Dự đoán lượng khách 4 quý tiếp theo}
\end{figure}

Dự đoán SARIMA cho 4 quý tiếp theo:

\begin{longtable}[]{@{}lll@{}}
\toprule\noalign{}
Quý & Dự đoán & Khoảng tin cậy 95\% \\
\midrule\noalign{}
\endhead
\bottomrule\noalign{}
\endlastfoot
2027 Q1 & 1,421,545 & {[}-247,180 --- 3,090,271{]} \\
2027 Q2 & -209,643 & {[}-3,479,627 --- 3,060,341{]} \\
2027 Q3 & -1,549,742 & {[}-5,812,211 --- 2,712,728{]} \\
2027 Q4 & -1,132,213 & {[}-6,199,123 --- 3,934,697{]} \\
\end{longtable}

\textbf{Lưu ý:} Khoảng tin cậy rất rộng, phản ánh mức độ không chắc chắn
cao do COVID-19 và dữ liệu hạn chế.

\newpage

\section{8. Tổng kết}\label{tux1ed5ng-kux1ebft}

\subsection{8.1 Kết quả đạt
được}\label{kux1ebft-quux1ea3-ux111ux1ea1t-ux111ux1b0ux1ee3c}

\begin{enumerate}
\def\labelenumi{\arabic{enumi}.}
\tightlist
\item
  Phân tích thành công dữ liệu du lịch 2008--2026,chỉ ra xu hướng tăng
  trưởng và tác động COVID-19.
\item
  Xác định Trung Quốc và Hàn Quốc là thị trường nguồn lớn nhất.
\item
  Phát hiện tính mùa vụ: Q1 cao nhất.
\item
  Xây dựng 5 mô hình, Chronos-T5-Base cho MAE thấp nhất.
\item
  Dự đoán 4 quý tiếp theo với khoảng tin cậy.
\end{enumerate}

\subsection{8.2 Hạn chế}\label{hux1ea1n-chux1ebf}

\begin{enumerate}
\def\labelenumi{\arabic{enumi}.}
\tightlist
\item
  Dữ liệu hạn chế (55 quý huấn luyện).
\item
  Khoảng trống 2021 phá vỡ tính liên tục chuỗi thời gian.
\item
  Thiếu đặc trưng bên ngoài (chính sách visa, tỷ giá\ldots).
\item
  Dự đoán SARIMA không ổn định.
\end{enumerate}

\subsection{8.3 Hướng phát
triển}\label{hux1b0ux1edbng-phuxe1t-triux1ec3n}

\begin{enumerate}
\def\labelenumi{\arabic{enumi}.}
\tightlist
\item
  Thu thập dữ liệu theo tháng để tăng mẫu.
\item
  Bổ sung đặc trưng bên ngoài (visa, tỷ giá, sự kiện du lịch).
\item
  Thử nghiệm Prophet, LSTM.
\item
  Dự đoán theo từng quốc gia riêng lẻ.
\item
  Xây dựng dashboard tương tác.
\end{enumerate}

\end{document}
